% Generated draft: requires mathematical and source review.
\documentclass[11pt]{article}
\usepackage[margin=25mm]{geometry}
\usepackage{fontspec}
\setmainfont{Latin Modern Roman}
\usepackage{amsmath,amssymb,mathtools,unicode-math}
\setmathfont{Latin Modern Math}
\usepackage{enumitem,needspace,xurl,hyperref}
\setlength{\parindent}{0pt}
\setlength{\parskip}{5pt}
\setlength{\emergencystretch}{3em}
\newcommand{\N}{\mathbb{N}}
\newcommand{\Z}{\mathbb{Z}}
\newcommand{\Q}{\mathbb{Q}}
\newcommand{\R}{\mathbb{R}}
\newcommand{\C}{\mathbb{C}}
\newcommand{\F}{\mathbb{F}}
\newcommand{\A}{\mathbb{A}}
\newcommand{\PP}{\mathbb{P}}
\newcommand{\E}{\mathbb{E}}
\newcommand{\eps}{\varepsilon}
\newcommand{\dd}{\,\mathrm{d}}
\begin{document}
\section*{UnsolvedMath ID: 800150}
\textbf{Source problem code:} catalogue entry
\subsection{Definitions and mathematical statement}
The supplied record contains two requests whose relationship is unclear.

For the Chebotarev request, a general framework is as follows. Let \(K\) be a number field, let \(L/K\) be a finite Galois extension with group \(G\), and let \(C_1,C_2\subseteq G\) be distinct nonempty conjugacy-invariant subsets. For \(x\ge 2\), define
\[
\pi_{C_i}(x;L/K)=\#\{\mathfrak p\subseteq\mathcal O_K:\mathfrak p\text{ is a nonzero prime ideal unramified in }L,\ N\mathfrak p\le x,\ \operatorname{Frob}_{\mathfrak p}\subseteq C_i\},
\]
where \(N\mathfrak p=|\mathcal O_K/\mathfrak p|\) and \(\operatorname{Frob}_{\mathfrak p}\) denotes the Frobenius conjugacy class. A two-way Chebotarev race compares these counting functions, possibly after normalization by \(|C_i|/|G|\). The supplied request is to investigate further examples and determine the most biased such races. However, the record specifies neither the family of extensions and subsets over which to optimize, nor the normalization, nor a numerical definition of bias. Thus this request does not yet determine a precise extremal problem. The framework above supplies terminology, not additional hypotheses attributed to the source.

The second passage, attributed to Greg Martin, reports an observation of Shanks from 1959. For a positive integer \(n\), let \(\Omega(n)\) be the number of prime factors of \(n\), counted with multiplicity, with \(\Omega(1)=0\). Define
\[
\chi_{-4}(n)=
\begin{cases}
0,&n\equiv0,2\pmod4,\\
1,&n\equiv1\pmod4,\\
-1,&n\equiv3\pmod4.
\end{cases}
\]
For each integer \(k\ge1\), each \(a\in\{1,3\}\), and each real \(x\ge1\), put
\[
A_{k,a}(x)=\#\{n\in\mathbb Z_{\ge1}:n\le x,\ \Omega(n)=k,\ n\equiv a\pmod4\},
\qquad
B_k(x)=(-1)^k\bigl(A_{k,1}(x)-A_{k,3}(x)\bigr).
\]
The reported proposed phenomenon is that, for every fixed integer \(k\ge1\), odd integers with \(\Omega(n)=k\) favor the residue class \((-1)^k\pmod4\); in this notation, the proposed preference is toward \(B_k(x)>0\). Equivalently at the level of individual integers, for odd \(n\) with \(\Omega(n)=k\),
\[
(-1)^{\Omega(n)}\chi_{-4}(n)=1
\quad\Longleftrightarrow\quad
n\equiv(-1)^k\pmod4.
\]
The record identifies \(k=1\) with Chebyshev's original observation. It does not specify the meaning of statistical bias, any density or averaging measure, or an eventual-inequality assertion. Consequently no such quantitative assertion is imposed here. The final source sentence also describes a preference of \((-1)^{\Omega(n)}\chi_{-4}(n)\) for \(1\) over \(-1\); whether this is intended separately for each fixed \(k\) or over all odd integers is unresolved.
\subsection{Short English statement}
Investigate the most biased two-way Chebotarev races, with the admissible races and the measure of bias still unspecified. A second supplied passage reports Shanks's proposed preference: for every fixed \(k\ge1\), odd integers having exactly \(k\) prime factors counted with multiplicity should favor \((-1)^k\pmod4\). The statistical meaning of this preference and the relationship between the two passages require clarification.
\subsection{Sources}
\begin{itemize}
\item Local UnsolvedMath registry: permanent ID \(800150\), exact source problem code \texttt{catalogue entry}, title ``Comparative Prime Number Theory --- Problem 19.'' The registry attributes the record to Hamieh, Kadiri, Martin, and Ng, with proposed year 2024, and labels it open. This is a supplied status label, not an independently verified mathematical assessment.
\item Hamieh, Kadiri, Martin, and Ng, \emph{Comparative Prime Number Theory Problem List} (2024), source item Problem (\#19), \url{https://arxiv.org/abs/2407.03530}. The supplied provenance records access on 2026-10-16 and extraction by \texttt{arxiv\_latex}; these are supplied provenance details, not a claim of independent source inspection.
\item The embedded passage is attributed to Greg Martin, University of British Columbia. It attributes the 1959 observation to Shanks under the reference label [Shanks], whose full bibliographic entry is not supplied, and identifies the case \(k=1\) with Chebyshev's original observation.
\end{itemize}
\end{document}
